\section{凸集}

\subsection{基本凸集}
\subsubsection{直线}
\begin{definition}
	设$x_1,x_2\in\mathbb{R}^{n}$且$x_1\ne x_2$。若$\theta\in[0,1]$，称所有满足形式：
	\begin{equation*}
		y=\theta x_1+(1-\theta)x_2
	\end{equation*}
	的点$y$构成的集合为\gls{LineSegment}。若$\theta\in\mathbb{R}$，称所有满足上述形式的点$y$构成的集合为\gls{Line}。
\end{definition}
\begin{note}
	上述形式等价于：
	\begin{equation*}
		y=x_2+\theta(x_1-x_2)
	\end{equation*}
	这给出了另一种直观：$y$是基$x_2$和被$\theta$缩放的方向$x_1-x_2$的和。
\end{note}
\subsubsection{凸集}
\begin{definition}
	设$E\subseteq\mathbb{R}^{n}$。若对于任意不同的$x_1,x_2\in E$，由$x_1,x_2$确定的线段都在$E$中，则称$E$是\gls{ConvexSet}。
\end{definition}
\begin{definition}
	设$n\in\mathbb{N}^+$，$\seq{\theta}{n}\geqslant0$且$\sum\limits_{i=1}^{n}\theta_i=1$，$\seq{x}{n}\in\mathbb{R}^{n}$。称：
	\begin{equation*}
		\theta_1x_1+\theta_2x_2+\cdots+\theta_nx_n
	\end{equation*}
	为$\seq{x}{n}$构成的一个\gls{ConvexCombination}。
\end{definition}
\begin{definition}
	设$E\subseteq\mathbb{R}^{n}$。定义：
	\begin{equation*}
		\left\{\sum_{i=1}^{n}\theta_ix_i:x_i\in E,\;\theta_i\geqslant0,\;\sum_{i=1}^{n}\theta_i=1,\;n\in\mathbb{N}^+\right\}
	\end{equation*}
	为$E$的\gls{ConvexHull}，记作$\operatorname{Conv}(E)$。
\end{definition}
\begin{figure}[H]
	\centering
	\begin{tikzpicture}[every node/.style={font=\small}]
		\begin{scope}[shift={(0,0)}]
			\draw[fill=blue!8, draw=blue!65!black, thick] (0,0) ellipse (1.65 and 1.15);
			\coordinate (a) at (-0.9,-0.35);
			\coordinate (b) at (0.95,0.45);
			\draw[very thick, teal!75!black] (a) -- (b);
			\fill[blue!75!black] (a) circle (2pt) node[below left] {$x_1$};
			\fill[blue!75!black] (b) circle (2pt) node[above right] {$x_2$};
			\node[font=\small\bfseries] at (0,-1.55) {凸集：线段始终留在集合内};
		\end{scope}
		\begin{scope}[shift={(5,0)}]
			\draw[fill=orange!8, draw=orange!85!black, thick]
			(-1.65,-1.05) -- (1.65,-1.05) -- (1.65,1.05) -- (0.45,1.05)
			-- (0.45,-0.1) -- (-0.45,-0.1) -- (-0.45,1.05) -- (-1.65,1.05) -- cycle;
			\coordinate (c) at (-1.05,0.65);
			\coordinate (d) at (1.05,0.65);
			\draw[very thick, red!70!black, dashed] (c) -- (d);
			\fill[orange!85!black] (c) circle (2pt) node[above] {$x_1$};
			\fill[orange!85!black] (d) circle (2pt) node[above] {$x_2$};
			\node[font=\small\bfseries] at (0,-1.55) {非凸集：部分线段离开集合};
		\end{scope}
	\end{tikzpicture}
	\caption{凸集与非凸集的二维几何结构。}
\end{figure}
\begin{property}\label{prop:ConvexSet}
	设$E\subseteq\mathbb{R}^{n}$。凸集具有如下性质：
	\begin{enumerate}
		\item $E$是凸集当且仅当$E$中任意有限个元素的凸组合都在$E$中；
		\item $E$的凸包是包含$E$的最小凸集；
		\item 凸集的交集是凸集；
		\item 设$A\in M_{m\times n}(\mathbb{R}),\;b\in\mathbb{R}^{n}$，$E$是凸集，则$AE+b=\{Ax+b:x\in E\}$也是凸集；
		\item 设$E_1,E_2$是凸集，则$E_1\pm E_1=\{x\pm y:x\in E_1,\;y\in E_2\}$也是凸集；
		\item 设$E_1,E_2$是凸集，则$\{(x,y_1+y_2):(x,y_1)\in E_1,\;(x,y_2)\in E_2\}$也是凸集；
		\item 若$E$是凸集，则$\{x_1\in\mathbb{R}^{m}:(x_1,x_2)\in E,\;m\in\mathbb{N}^+,\;m\leqslant n\}$也是凸集；
		\item 若$E$是凸集，则$\overline{E}$也是凸集；
	\end{enumerate}
\end{property}
\begin{proof}
	(1)\textbf{必要性：}当$n=2$时，由凸集的定义可知$E$中任意$2$个元素的凸组合都在$E$中。假设结论对$n-1$成立，下面证明对$n$个也成立。\par
	任取$\seq{\theta}{n}\geqslant0$且$\sum\limits_{i=1}^{n}\theta_i=1$，$\seq{x}{n}\in E$，则：
	\begin{equation*}
		\sum_{i=1}^{n}\theta_ix_i=\sum_{i=1}^{n-1}\theta_ix_i+\theta_nx_n=\sum_{i=1}^{n-1}\theta_ix_i+\left(1-\sum_{i=1}^{n-1}\theta_i\right)x_i
	\end{equation*}
	若$\sum\limits_{i=1}^{n-1}\theta_i=0$，则上式退化为$x_n$，结论成立。若$\sum\limits_{i=1}^{n-1}\theta_i\ne0$，则：
	\begin{equation*}
		y=\sum_{i=1}^{n-1}\frac{\theta_i}{\sum\limits_{j=1}^{n-1}\theta_j}x_i
	\end{equation*}
	是关于$\seq{x}{n-1}$的一个凸组合，由归纳假设可知$y\in E$。变形可得：
	\begin{equation*}
		\sum_{i=1}^{n}\theta_ix_i=\sum_{i=1}^{n-1}\theta_ix_i+\left(1-\sum_{i=1}^{n-1}\theta_i\right)x_i=\sum\limits_{j=1}^{n-1}\theta_jy+\left(1-\sum_{i=1}^{n-1}\theta_i\right)x_i\in E
	\end{equation*}
	必要性得证。\par
	\textbf{充分性：}任意有限个包含$2$个。\par
	(2)由(1)立即可得。\par
	(3)任取指标集$I$和凸集$C_i,i\in I$，令$C=\underset{i\in I}{\overset{}{\cap}}C_i$。任取$x_1,x_2\in C$和$\theta\geqslant0$，则对任意的$i\in I$有$x_1,x_2\in C_i,\;\theta x_1+(1-\theta)x_2\in C_i$，于是$\theta x_1+(1-\theta)x_2\in C$，$C$是凸集。\par
	(4)(5)(6)(7)利用凸集的定义可直接得到证明，略去。\par
	(8)任取$x,y\in\overline{E}$，于是存在$E$中的$\{x_n\}$和$\{y_n\}$满足$\{x_n\}\to x$和$\{y_n\}\to y$。任取$\theta\in[0,1]$，因为$E$是凸集，所以$\theta x_n+(1-\theta)y_n\in E$。由\cref{prop:Norm}(3)可知$\theta x+(1-\theta)y\in\overline{E}$。根据$x,y$的任意性可得$\overline{E}$是凸集。
\end{proof}
\subsubsection{仿射集}
\begin{definition}
	设$E\subseteq\mathbb{R}^{n}$。若对于任意不同的$x_1,x_2\in E$，由$x_1,x_2$确定的直线都在$E$中，则称$E$是\gls{AffineSet}。
\end{definition}
\begin{definition}
	设$n\in\mathbb{N}^+$，$\seq{\theta}{n}\in\mathbb{R}$且$\sum\limits_{i=1}^{n}\theta_i=1$，$\seq{x}{n}\in\mathbb{R}^{n}$。称：
	\begin{equation*}
		\theta_1x_1+\theta_2x_2+\cdots+\theta_nx_n
	\end{equation*}
	为$\seq{x}{n}$构成的一个\gls{AffineCombination}。
\end{definition}
\begin{definition}
	设$E\subseteq\mathbb{R}^{n}$。定义：
	\begin{equation*}
		\left\{\sum_{i=1}^{n}\theta_ix_i:x_i\in E,\;\theta_i\in\mathbb{R},\;\sum_{i=1}^{n}\theta_i=1,\;n\in\mathbb{N}^+\right\}
	\end{equation*}
	为$E$的\gls{AffineHull}，记作$\operatorname{Aff}(E)$。
\end{definition}
\begin{property}\label{prop:AffineSet}
	设$E\subseteq\mathbb{R}^{n}$。仿射集具有如下性质：
	\begin{enumerate}
		\item $E$是仿射集当且仅当$E$中任意有限个元素的仿射组合都在$E$中；
		\item $E$是仿射集当且仅当存在一个$\mathbb{R}^{n}$的子空间$V$使得$E=\{x+x_0:x\in V,\;x_0\in E\}$，并且表达式与$x_0$的选择无关；
		\item $E$的仿射包是包含$E$的最小仿射集；
		\item 仿射集是凸集；
	\end{enumerate}
\end{property}
\begin{proof}
	(1)完全类似于\cref{prop:ConvexSet}(1)。\par
	(2)\textbf{必要性：}设$E$是仿射集，任取$x_0\in E$，令$V=\{x-x_0:x\in E\}$，则$E=\{x+x_0:x\in V,\;x_0\in E\}$。\par
	任取$\alpha,\beta\in\mathbb{R}$和$x_1,x_2\in V$，则：
	\begin{equation*}
		\alpha x_1+\beta x_2+x_0=\alpha(x_1+x_0)+\beta(x_2+x_0)+(1-\alpha-\beta)x_0\in E
	\end{equation*}
	于是$\alpha x_1+\beta x_2\in V$。由\cref{theo:Subspace}可知$V$是子空间。\par
	取不同于$x_0$的$x_1\in E$构成集合$V'=\{x-x_1:x\in C\}$。任取$V'$中的元素$\alpha=x-x_1$，则$\alpha=x+x_1-x_0$是$x,x_1,x_0$构成的一个仿射组合，于是$\alpha\in V$。由$\alpha$的任意性可知$V'\subseteq V$。同理可得$V\subseteq V'$，于是有$V=V'$，即表达式与$x_0$的选择无关。\par
	\textbf{充分性：}任取$E$中不同的两点$x_1+x_0,x_2+x_0$和$\theta\in\mathbb{R}$，有：
	\begin{equation*}
		\theta(x_1+x_0)+(1-\theta)(x_2+x_0)=\theta x_1+(1-\theta)x_2+x_0
	\end{equation*}
	因为$V$是子空间，由\cref{theo:Subspace}可知$\theta x_1+(1-\theta)x_2\in V$，所以$\theta x_1+(1-\theta)x_2+x_0\in E$。由$x_1+x_0,x_2+x_0$和$\theta$的任意性可知$E$是仿射集。\par
	(3)由(1)立即可得。\par
	(4)由定义立即可得。
\end{proof}
\begin{definition}
	设$E\subseteq\mathbb{R}^{n}$是仿射集。由\cref{prop:AffineSet}(2)可知存在一个$\mathbb{R}^{n}$的子空间$V$使得$E=\{x+x_0:x\in V,\;x_0\in E\}$，称$V$是$E$的方向子空间，记$\dim E=\dim V$。
\end{definition}
\begin{definition}
	设$x_0,x_1,\dots,x_m\in\mathbb{R}^n$。若：
	\begin{equation*}
		\sum_{i=0}^{m}\lambda_ix_i=0,\quad\sum_{i=0}^{m}\lambda_i=0
	\end{equation*}
	只能推出：
	\begin{equation*}
		\lambda_0=\lambda_1=\cdots=\lambda_m=0
	\end{equation*}
	则称$x_0,x_1,\dots,x_m$\gls{AffineIndependent}。
\end{definition}
\begin{property}\label{prop:AffineIndependence}
	设$x_0,x_1,\dots,x_m\in\mathbb{R}^n$。仿射无关具有以下性质：
	\begin{enumerate}
		\item $x_0,x_1,\dots,x_m$仿射无关当且仅当对任意的或某个$j\in\{0,1,\dots,m\}$有$\{x_i-x_j:i\neq j\}$线性无关；
		\item 对任意的$a\in\mathbb{R}^n$，$x_0,x_1,\dots,x_m$仿射无关当且仅当$x_0+a,x_1+a,\dots,x_m+a$仿射无关；
		\item 若$x_0,x_1,\dots,x_m$仿射无关，则对于任意的$j\in\{0,1,\dots,m\}$有$\operatorname{Aff}(\{x_0,x_1,\dots,x_m\})=x_j+\operatorname{span}\{x_i-x_j:i\ne j\}$且$\dim\operatorname{Aff}(\{x_0,x_1,\dots,x_m\})=m$；
		\item 若$E$是维数为$m$的仿射集，则存在$m+1$个仿射无关的点$x_0,x_1,\dots,x_m\in E$使得$E=\operatorname{Aff}(\{x_0,x_1,\dots,x_m\})$；
	\end{enumerate}
\end{property}
\begin{proof}
	(1)\textbf{必要性：}任取$j\in\{0,1,\dots,m\}$，设：
	\begin{equation*}
		\sum_{i\ne j}\lambda_i(x_i-x_j)=0
	\end{equation*}
	于是：
	\begin{equation*}
		\sum_{i\ne j}\lambda_i(x_i-x_j)=\sum_{i\ne j}\lambda_ix_i-\sum_{i\ne j}\lambda_ix_j
	\end{equation*}
	记$\lambda_j=-\sum\limits_{i\ne j}\lambda_i$，则；
	\begin{equation*}
		\sum_{i\ne j}\lambda_i(x_i-x_j)=\sum_{i=0}^{m}\lambda_ix_i,\quad\sum_{i=0}^{m}\lambda_i=0
	\end{equation*}
	因为$x_0,x_1,\dots,x_m$仿射无关，所以对任意的$i\ne j$都有$\lambda_i=0$，即$\{x_i-x_j:i\neq j\}$线性无关。由$j$的任意性可知必要性成立。\par
	\textbf{充分性：}任取$j\in\{0,1,\dots,m\}$，设：
	\begin{equation*}
		\sum_{i=0}^{m}\lambda_ix_i,\quad\sum_{i=0}^{m}\lambda_i=0
	\end{equation*}
	则$\lambda_j=-\sum\limits_{i\ne j}\lambda_i$，代入上式可得：
	\begin{equation*}
		\sum_{i=0}^{m}\lambda_ix_i=\sum_{i\ne j}\lambda_ix_i-\sum_{i\ne j}\lambda_ix_j=\sum_{i\ne j}^{}\lambda_i(x_i-x_j)
	\end{equation*}
	因为$\{x_i-x_j:i\ne j\}$线性无关，所以对任意的$i\ne j$都有$\lambda_i=0$，于是$\lambda_j=0$，即$x_0,x_1,\dots,x_m$仿射无关。\par
	(2)\textbf{必要性：}若：
	\begin{equation*}
		\sum_{i=0}^{m}\lambda_i(x_i+a)=0,\quad\sum_{i=0}^{m}\lambda_i=0
	\end{equation*}
	注意到：
	\begin{equation*}
		\sum_{i=0}^{m}\lambda_i(x_i+a)=\sum_{i=0}^{m}\lambda_ix_i+\sum_{i=0}^{m}\lambda_ia=\sum_{i=0}^{m}\lambda_ix_i
	\end{equation*}
	因为$x_0,x_1,\dots,x_m$仿射无关，所以$\lambda_0=\lambda_1=\cdots=\lambda_m=0$，于是$x_0+a,x_1+a,\dots,x_m+a$仿射无关。\par
	\textbf{充分性：}由定义立即可得。\par
	(3)任取$x\in\operatorname{Aff}(\{x_0,x_1,\dots,x_m\})$和$j\in\{0,1,\dots,m\}$，则存在$\lambda_i\in\mathbb{R},\;i\ne j$使得：
	\begin{equation*}
		x=\sum_{i\ne j}\lambda_ix_i+\left(1-\sum_{i\ne j}\lambda_i\right)x_j=x_j+\sum_{i\ne j}\lambda_i(x_i-x_j)
	\end{equation*}
	由$x$和$j$的任意性可知对于任意的$j\in\{0,1,\dots,m\}$有$\operatorname{Aff}(\{x_0,x_1,\dots,x_m\})=x_j+\operatorname{span}\{x_i-x_j:i\ne j\}$。\par
	根据\cref{prop:SpanSubspace}(1)、\cref{prop:AffineSet}(2)、(1)和\cref{prop:SpanSubspace}(2)可知$\dim\operatorname{Aff}(\{x_0,x_1,\dots,x_m\})=m$。\par
	(4)任取$x_0\in E$，由\cref{prop:AffineSet}(2)可知存在一个$\mathbb{R}^{n}$的子空间$V$使得$E=\{x+x_0:x\in V\}$。取$V$的一组基$\seq{\alpha}{m}$，记$x_i=x_0+v_i\in E,\;i=1,2,\dots,m$。根据(1)可知$x_0,x_1,\dots,x_m$仿射无关。由(3)和\cref{prop:AffineSet}(2)可知结论成立。
\end{proof}
\subsubsection{相对内部}
\begin{definition}
	设$E$是$\mathbb{R}^{n}$的非空子集，$A=\operatorname{Aff}(E)$。定义：
	\begin{equation*}
		\operatorname{ri}(E)\coloneq\{x\in E:\exists\;r>0,\;U(x,r)\cap A\nsubseteq E\}
	\end{equation*}
	为$E$的\gls{RelativeInterior}，其中的点称为$E$的\textbf{相对内点}。约定$\operatorname{ri}(\varnothing)=\varnothing$。
\end{definition}
\begin{property}\label{prop:RelativeInterior}
	设$E\subseteq\mathbb{R}^{n}$为非空凸集，则：
	\begin{enumerate}
		\item $\operatorname{ri}(E)\ne\varnothing$；
		\item 若$x\in\operatorname{ri}(E),\;y\in\overline{E}$，则$(1-\theta)x+\theta y\in\operatorname{ri}(E)$，$\forall\;\theta\in[0,1)$；特别地，$\operatorname{ri}(E)$是凸集；
		\item $\operatorname{ri}(\overline{E})=\operatorname{ri}(E)$；
	\end{enumerate}
\end{property}
\begin{proof}
	(1)设$A=\operatorname{Aff}(E)$，根据\cref{prop:AffineSet}(2)可知存在一个$\mathbb{R}^{n}$的子空间$V$使得$A=\{x+x_0:x\in V,x_0\in E\}$，记$\dim E=\dim V=m$。\par
	若$m=0$，则$E$为单点集，结论成立。\par
	若$m>0$，
\end{proof}
\subsection{分离超平面定理}
\begin{definition}
	设$a\in\mathbb{R}^{n},\;b\in\mathbb{R}$。称$\{x\in\mathbb{R}^{n}:a^{\top}x=b\}$为\gls{Hyperplane}，$\{x\in\mathbb{R}^{n}:a^{\top}x\leqslant b\}$为闭\gls{Halfspace}，$\{x\in\mathbb{R}^{n}:a^{\top}x< b\}$为开半空间。
\end{definition}
\begin{lemma}\label{lem:BestApproximationElement}
	设$X$是实Hilbert空间，$E$是$X$中的凸闭集。对任意的$x\in X$，$y$是$x$在$E$中的最佳逼近元当且仅当：
	\begin{equation*}
		\forall\;z\in E,\;(x-y,z-y)\leqslant0
	\end{equation*}
\end{lemma}
\begin{proof}
	\textbf{(1)必要性：}对任意的$z\in E$和$\theta\in[0,1]$，因为$E$是$X$中的凸集，所以$\theta z+(1-\theta)y\in E$。由于$y$是$x$在$E$中的最佳逼近元，结合：
	\begin{equation*}
		||x-y||^2\leqslant||x-[\theta z+(1-\theta)y]||^2=||x-y-\theta(z-y)||^2=||x-y||^2-2\theta(x-y,z-y)+\theta^2||z-y||^2
	\end{equation*}
	可知：
	\begin{equation*}
		-2\theta(x-y,z-y)+\theta^2||z-y||^2\geqslant0
	\end{equation*}
	取$\theta\to0^+$，由\cref{prop:RMap}(4)可知$(x-y,z-y)\leqslant0$。\par
	\textbf{(2)充分性：}此时对任意的$z\in E$有：
	\begin{align*}
		&||x-z||=||x-y+y-z||^2=||x-y||^2+2(x-y,y-z)+||y-z||^2 \\
		=&|x-y||^2-2(x-y,z-y)+||y-z||^2\geqslant||x-y||
	\end{align*}
	于是$y$是$x$在$E$中的最佳逼近元。
\end{proof}
\begin{theorem}[Hyperplane Separation Theorem]\label{theo:HyperplaneSeparationTheorem}
	设$A$和$B$是$\mathbb{R}^{n}$上不交的非空凸集，则存在非零的$a\in\mathbb{R}^{n}$和$b\in\mathbb{R}$使得：
	\begin{equation*}
		\forall\;x\in A,\;(a,x)\geqslant b,\quad\forall\;y\in B,\;(a,y)\leqslant b
	\end{equation*}
	记$C=A-B=\{x-y:x\in A,\;y\in B\}$，若$\mathbf{0}\notin\overline{C}$，则存在$a\in\mathbb{R}^{n}$和$b\in\mathbb{R}$使得上式取严格的不等号。
\end{theorem}
\begin{proof}
	由\cref{prop:ConvexSet}(5)(8)可知$C$是凸集且$\overline{C}$是闭凸集。\par
	若$\mathbf{0}\notin\overline{C}$，由\cref{theo:BestApproximationElement}可知在$\overline{C}$中存在$\mathbf{0}$的唯一的最佳逼近元$a$。根据\cref{lem:BestApproximationElement}可得：
	\begin{equation*}
		\forall\;z\in\overline{C},\;(-a,z-a)\leqslant0
	\end{equation*}
	也即$(a,z)\geqslant||a||^2$。对任意的$x\in A$和$y\in B$有$x-y\in\overline{C}$，所以：
	\begin{equation*}
		(a,x)\geqslant(a,y)+||a||^2
	\end{equation*}
	根据上式与\info{确界原理}，可记：
	\begin{equation*}
		\alpha=\inf_{x\in A}(a,x),\quad\beta=\sup_{y\in B}(a,y)
	\end{equation*}
	由上下确界的不等式性可知：
	\begin{equation*}
		\alpha\geqslant\beta+||a||^2>\beta
	\end{equation*}
	任取$b\in(\beta,\alpha)$可得：
	\begin{equation*}
		\forall\;x\in A,\;(a,x)\geqslant\alpha>b,\quad\forall\;y\in B,\;(a,y)\leqslant\beta<b
	\end{equation*}\par
	若$\mathbf{0}\in\overline{C}$，因为$\mathbf{0}\not\in C$且$\operatorname{ri}(C)\subseteq C$，所以$\mathbf{0}\notin\operatorname{ri}(C)$。由\cref{prop:RelativeInterior}(3)可知$\mathbf{0}\notin\operatorname{ri}(\overline{C})$，即：
	\begin{equation*}
		\forall\;r>0,\;U(\mathbf{0},r)\cap\operatorname{Aff}(\overline{C})\nsubseteq\overline{C}
	\end{equation*}
	于是对任意的$n\in\mathbb{N}^+$，存在$x_n\in\operatorname{Aff}(\overline{C})\backslash\overline{C}$满足$||x_n||\leqslant\dfrac{1}{n}$。根据\cref{theo:BestApproximationElement}可知在$\overline{C}$中存在$x_n$的唯一的最佳逼近元$y_n$。由\cref{lem:BestApproximationElement}可得：
	\begin{equation*}
		\forall\;z\in\overline{C},\;(x_n-y_n,z-y_n)\leqslant0
	\end{equation*}
	即：
	\begin{equation*}
		\forall\;z\in\overline{C},\;\left(\frac{y_n-x_n}{||y_n-x_n||},z\right)\geqslant\left(\frac{y_n-x_n}{||y_n-x_n||},y_n\right)
	\end{equation*}
	因为单位球面是有界闭集（因为$\{1\}$是闭集，根据\cref{prop:Norm}(2)和\cref{theo:ContinousMapO2OC2C}可知$\{x;||x||_2=1\}$是闭集。），由\cref{theo:CompactRn}可知$\left\{\dfrac{y_n-x_n}{||y_n-x_n||}\right\}\to a$且$||a||=1$。根据最佳逼近元的定义可知：
	\begin{equation*}
		||y_n||=||y_n-x_n+x_n-\mathbf{0}||\leqslant||x_n-y_n||+||x_n||\leqslant||x_n-\mathbf{0}||+||x_n||\to0
	\end{equation*}
	所以$\{y_n\}\to\mathbf{0}$。由\cref{prop:InnerProduct}(2)(1)可知：
	\begin{equation*}
		\forall\;z\in\overline{C},\;(a,z)\geqslant0
	\end{equation*}
	对任意的$x\in A$和$y\in B$有$x-y\in\overline{C}$，所以：
	\begin{equation*}
		(a,x)\geqslant(a,y)
	\end{equation*}
	根据上式与\info{确界原理}，可记：
	\begin{equation*}
		\alpha=\inf_{x\in A}(a,x),\quad\beta=\sup_{y\in B}(a,y)
	\end{equation*}
	由上下确界的不等式性可知$\alpha\geqslant\beta$。任取$b\in[\beta,\alpha]$可得：
	\begin{equation*}
		\forall\;x\in A,\;(a,x)\geqslant b,\quad\forall\;y\in B,\;(a,y)\leqslant b\qedhere
	\end{equation*}
\end{proof}
\begin{definition}
	设$E\subseteq\mathbb{R}^{n}$，$x_0\in\partial E$。若非零的$\alpha\in\mathbb{R}^{n}$和$b\in\mathbb{R}$满足对任意的$x\in E$有$(\alpha,x)\leqslant(\alpha,x_0)$，则称超平面$\{x\in\mathbb{R}^{n}:(\alpha,x)=(\alpha,x_0)\}$为$E$在$x_0$处的\gls{SupportingHyperplane}。
\end{definition}
\begin{theorem}[Supporting Hyperplane Theorem]\label{theo:SupportingHyperplaneTheorem}
	若$E\subseteq\mathbb{R}^{n}$是凸集，则在$E$的任意边界点处都存在支撑超平面。
\end{theorem}
